\section{JAX versus PyTorch Benchmark}
\label{sec:jax-vs-pytorch}

\subsection{Architecture matching}

The benchmark implementation provides a PyTorch counterpart to the trained
Flax residual network.  Both use three encoder levels with 16, 32, and 64
channels; paired \(3\times3\) convolutions; GELU activation; stride-two
downsampling; transposed-convolution upsampling; skip concatenation; and a
zero-initialised \(1\times1\), three-channel output head.  Flax uses NHWC
layout and PyTorch uses NCHW layout.  Layout conversion and host-to-device
transfer are deliberately placed outside timed regions.  Automated
architecture checks report 150,883 trainable scalar parameters for each
implementation.

\subsection{Benchmark protocol}

The publication configuration specifies identical seeded float32 inputs with
six channels and a spatial shape of 128 by 128 at batch sizes 1, 4, 8, and
16.  For each batch size and framework, ten warm-up iterations precede 50
measured iterations.  Forward inference and a forward--backward AdamW update
are timed separately.  JAX compilation and PyTorch's compiled first iteration
are recorded separately from steady-state measurements.  Every GPU interval
is synchronised before the wall-clock measurement is closed; the output schema
contains the sample mean and sample standard deviation rather than a single
run.

Each batch size runs in a fresh worker process.  This design avoids carrying
JAX's process-level allocator peak from one batch into the next, while
PyTorch's peak allocation is reset after warm-up.  JAX profiler traces and
PyTorch Chrome traces are aggregated into a common operation-duration table.
The configuration requires one NVIDIA GPU to be simultaneously visible to
both frameworks and explicitly forbids automatic CPU substitution for
publication results.

\subsection{Execution status}

The final benchmark was not executable on the present host.  JAX exposed only
\texttt{CpuDevice(id=0)}, PyTorch reported
\texttt{torch.cuda.is\_available()=False}, and the runtime probe found no
NVIDIA CUDA driver library.  The benchmark therefore terminated before
measuring batch size 1, as designed.  No results CSV, LaTeX table, profiling
summary, or benchmark bar chart was generated.

Consequently, this report makes no empirical claim about JAX-versus-PyTorch
GPU speed, compilation cost, or memory use.  Completing this chapter requires
running \texttt{scripts/08\_benchmark\_jax\_vs\_pytorch.py} with the locked
environment on a host where both frameworks see the same NVIDIA GPU.  CPU
smoke-test timings, even if produced, would not answer the stated GPU research
question and are not substituted here.
